\documentclass[11pt]{article}

\usepackage[margin=1in]{geometry}
\usepackage{amsmath,amssymb,bm,booktabs,microtype}
\usepackage[hidelinks]{hyperref}
\usepackage[nameinlink,noabbrev]{cleveref}
\usepackage[round,authoryear]{natbib}

\newcommand{\dd}{\,\mathrm{d}}
\newcommand{\E}{\mathbb{E}}
\newcommand{\Var}{\operatorname{Var}}
\newcommand{\taugeo}{\tau_{\mathrm{geo}}}
\newcommand{\lp}{\ell'}

\title{From the Equivalence Theorem to Photon-Path Cumulant Fitting\\
\large Derivation, assumptions, literature context, and application to OCO-2 spectra}
\author{OCO-2 cloud-proximity analysis: technical report}
\date{18 July 2026}

\begin{document}
\maketitle

\begin{abstract}
This report derives the spectral transmittance fit used in the OCO-2
cloud-proximity analysis from the radiative-transfer equivalence theorem.  In a
narrow band such as the oxygen A band, the continuum scattering properties are
assumed to be effectively wavelength independent.  A single conservative
scattering calculation then defines an ensemble of photon trajectories, and
gas absorption reweights each trajectory by its Beer--Lambert survival
probability.  Under the additional scalar-path approximation, the normalized
radiance is the Laplace transform of a wavelength-independent photon
path-length distribution.  Expanding its logarithm gives the fitted mean path
enhancement, path variance, and higher scaled cumulants.  We state the
assumption required at every step, distinguish the underlying path ensemble
from the absorption-selected paths actually detected in strong and weak
channels, and explain how the classical literature maps onto the equations and
the current implementation.
\end{abstract}

\section{Purpose and scope}

The practical question is whether a resolved absorption spectrum can be used
to infer information about the photon paths created by clouds, aerosols, the
surface, and three-dimensional transport.  The foundational result is the
equivalence theorem of \citet{irvine1964}, developed for high-resolution
atmospheric calculations by \citet{partain2000} and applied to molecular-line
absorption and the O$_2$ A band by \citet{stephens2000} and
\citet{heidinger2000}.  High-resolution measurements have subsequently been
used to infer photon-path distributions or their moments
\citep{funk2003,min2004,scholl2006}.

The present project does not invert a complete path-length probability density
function (PDF).  It fits a truncated expansion of log transmittance against
modeled slant optical depth, separately for O$_2$ A, weak CO$_2$, and strong
CO$_2$.  The goal of this report is to show exactly when the fitted
coefficients have a photon-path interpretation and when they should instead be
called effective spectral-shape coefficients.  This report is the expanded
theoretical companion to \texttt{FITTING\_DERIVATION.md}; both documents follow
the same project sequence: conservative trajectory ensemble, Beer--Lambert
weighting, scalar-path reduction, Laplace transform, continuum normalization,
log-cumulant expansion, and exact linear fitting.

\section{Objects that must be kept distinct}

Let $\Gamma$ denote a complete photon trajectory from solar illumination to the
detector, including every scattering and surface interaction.  Three related
distributions appear in the argument:

\begin{enumerate}
  \item $P_0(\Gamma)$ is the trajectory distribution in the corresponding
        conservative-scattering problem, before the target gas absorption is
        applied.
  \item $p_0(\lp)$ is the distribution of a scalar normalized path coordinate
        $\lp$ induced by $P_0(\Gamma)$.
  \item $p_{\mathrm{det}}(\lp\mid\tau)$ is the distribution among photons that
        survive absorption and reach the detector at optical depth $\tau$.
\end{enumerate}

The equivalence theorem requires the first distribution to remain unchanged
when absorption is introduced.  It does \emph{not} require the detected-path
distribution to be the same in weak and strong absorption channels.  Strong
absorption preferentially removes long paths, so a change in
$p_{\mathrm{det}}(\lp\mid\tau)$ is an expected prediction of the theorem, not
evidence against it.

\section{Equivalence theorem at the trajectory level}

Consider a narrow spectral interval and a fixed illumination--view geometry.
Let $I_{0,\lambda}$ denote the radiance for an otherwise equivalent reference
problem in which the target molecular absorption is absent.  The absorption
optical depth accumulated along trajectory $\Gamma$ is
\begin{equation}
  A_\lambda(\Gamma)
  = \int_\Gamma \alpha_\lambda(\bm r)\dd s,
  \label{eq:path_absorption}
\end{equation}
where $\alpha_\lambda$ is the local molecular absorption coefficient.  The
Beer--Lambert survival probability of that trajectory is
$\exp[-A_\lambda(\Gamma)]$.  Averaging over the conservative trajectory
ensemble gives
\begin{equation}
  \frac{I_\lambda}{I_{0,\lambda}}
  = \int P_0(\Gamma)
      \exp[-A_\lambda(\Gamma)]\dd\Gamma
  = \E_0\!\left\{\exp[-A_\lambda(\Gamma)]\right\}.
  \label{eq:general_equivalence}
\end{equation}

Equation~\eqref{eq:general_equivalence} is the general equivalence-theorem
statement relevant here.  It is independent of whether the scattering medium
is plane parallel or three dimensional.

\citet{stephens2000} formulated the O$_2$ A-band problem in terms of the
distribution of continuum photon optical paths.  Their equivalence-theorem
construction separates the continuum multiple-scattering calculation from
line absorption: a Monte Carlo calculation at a nonabsorbing wavelength first
provides the path distribution and continuum radiance, after which the
transmission function associated with gaseous absorption is integrated over
that distribution.  The same stored path statistics can therefore generate a
complete high-resolution line spectrum as the line-to-continuum absorption
ratio is varied.  They demonstrated this numerically by comparing spectra
reconstructed from the path distribution with direct multiple-scattering
calculations, finding essentially identical results for their test case.  They
also applied the construction grid point by grid point to horizontally
heterogeneous cloud fields: three-dimensional transport changes the path
distribution, but does not alter the equivalence-theorem weighting once that
distribution is known.  This is the precise sense in which
\cref{eq:general_equivalence} is geometry independent.  Their subsequent O$_2$
A-band application showed that the line-shape response carries information on
cloud optical depth, phase function, cloud pressure and pressure thickness,
and surface albedo \citep{heidinger2000}.  For this project, the key inheritance
is narrower: cloud geometry and 3-D transport are encoded in $P_0(\Gamma)$,
while wavelength-dependent O$_2$ absorption interrogates that ensemble.

\citet{partain2000} developed generalized, computationally useful forms of the
same theorem for line-by-line atmospheric radiative transfer.  Rather than
repeating a multiple-scattering solution at every spectral point, their method
accumulates probability distributions of photon histories during a reference
scattering calculation and applies spectral gas, particle, and surface
absorption afterward.  Use of geometrical path information permits different
vertical gas profiles and multiple homogeneous layers to be treated, while
joint path/scattering statistics extend the bookkeeping to nonconservative
particles and surface absorption.  The complete joint distribution is
memory-intensive, so they also introduced a reduced PDF approximation; later
analysis noted that compressing the full trajectory information can introduce
absorption bias.  This distinction is directly relevant here.  The present
project does not retain the complete layer-resolved or joint PDF used by the
generalized theorem.  It makes the stronger reduction
$A_\lambda(\Gamma)\simeq\tau_\lambda\lp(\Gamma)$ and fits a one-dimensional
marginal distribution.  Partain et al. therefore support the trajectory-level
starting point, while also showing what information is discarded by the
project's scalar-path approximation.

\paragraph{Assumptions at this step.}
\begin{enumerate}
  \item Introducing target-gas absorption does not change the geometry or the
        continuum scattering phase function, extinction, or boundary
        conditions.
  \item Inelastic processes that change photon wavelength, such as Raman
        scattering, are negligible or separately corrected.
  \item The reference radiance $I_{0,\lambda}$ represents the same scattering
        problem without the target absorption.  A continuum estimate that
        contains unmodeled surface or instrumental spectral structure violates
        this identification.
\end{enumerate}

\section{Narrow-band reduction to one path-length PDF}

For a narrow band such as O$_2$ A, cloud and aerosol scattering properties are
often assumed to vary slowly compared with the many-orders-of-magnitude change
in molecular line absorption.  This is the central spectral-separation
assumption used in O$_2$ A-band PPDF work
\citep{stephens2000,min2004,scholl2006}.

The three studies play complementary roles in the present derivation.
\citet{stephens2000} provide the theoretical separation: if continuum particle
properties and geometry are unchanged across the line, changes in spectral
radiance can be generated by applying the molecular transmission function to
one continuum path distribution.  Their path coordinate is naturally an
optical path of the continuum scattering problem.  Reorganizing a simulated
A-band spectrum by O$_2$ optical depth makes the line-growth behavior and its
dependence on the path distribution explicit.  In the project equations, this
motivates holding $P_0(\Gamma)$ fixed while scanning $\tau_\lambda$; it does not
by itself justify replacing the full trajectory-dependent absorber amount by
the scalar product $\tau_\lambda\lp$.

\citet{min2004} connect that theory to high-resolution observations.  Their
instrument resolved the O$_2$ A band and a nearby water-vapor band with better
than approximately $0.5\ \mathrm{cm}^{-1}$ resolution and out-of-band rejection
better than $10^{-5}$, properties needed because line-center leakage and the
instrument slit function strongly limit inverse-Laplace information.  They
wrote the multiple-path Beer--Lambert relation as an integral over the absorber
amount encountered along photon trajectories and treated the resolved
absorption line as an inverse-Laplace problem.  O$_2$ is advantageous because
its abundance profile is well known and its A-band lines span a large range of
opacity.  Their comparison with the nearby water-vapor band is especially
important for interpretation: the inferred mean and variance differed between
the two absorbers because water vapor is vertically inhomogeneous and coupled
to cloud location.  Thus a common geometrical scattering ensemble need not
produce the same \emph{effective absorber-path distribution} for different
gases.  For this project, that result supports separate O$_2$, weak-CO$_2$, and
strong-CO$_2$ fits and motivates calling their coefficients band-effective
absorber-path cumulants rather than universal geometrical-path moments.

\citet{scholl2006} make the vertical structure still more explicit.  Their
forward model used pressure- and temperature-dependent O$_2$ line absorption
and partitioned the atmospheric path into regions above the cloud, within the
cloud, and below the cloud, including an approximate cloud--surface
multiple-reflection contribution.  The random in-cloud path was represented by
a parameterized PDF and integrated against the wavelength-dependent absorption
of that region; the resulting high-resolution spectrum was convolved with the
measured instrumental slit function.  Although their spectral resolution and
out-of-band rejection supported roughly four independent pieces of
information, their primary analysis retained the first two path moments and
usually adopted a gamma-type PDF.  For sufficiently simple, well-layered
clouds and high signal-to-noise spectra, their sensitivity analysis indicated
that these moments could be recovered to approximately five percent.  Their
measurements were transmitted skylight rather than OCO-2 reflected radiance,
so their path geometry is not transferred directly to this project.  Their
method nevertheless identifies the exact approximation made here: the project
collapses their separate above-, in-, and below-cloud absorber paths into one
normalized multiplier $\lp$.  Agreement of the fitted OCO-2 coefficients with
layer-resolved Monte Carlo moments would therefore be a closure test of the
scalar reduction, not of the equivalence theorem itself.

To reduce \cref{eq:general_equivalence} to a one-dimensional Laplace transform,
make the stronger scalar-path approximation
\begin{equation}
  A_\lambda(\Gamma)
  \simeq \taugeo(\lambda)\,\lp(\Gamma).
  \label{eq:scalar_path}
\end{equation}
Here $\taugeo(\lambda)$ is the molecular optical depth along the direct
geometrical slant path, and $\lp$ is the photon absorption-path enhancement
relative to that direct path.  Defining
\begin{equation}
  p_0(\lp)
  = \int P_0(\Gamma)
      \delta\!\left[\lp-\lp(\Gamma)\right]\dd\Gamma,
  \qquad
  \int_0^\infty p_0(\lp)\dd\lp=1,
  \label{eq:induced_pdf}
\end{equation}
and substituting \cref{eq:scalar_path,eq:induced_pdf} into
\cref{eq:general_equivalence} yields
\begin{equation}
  \mathcal T(\taugeo)
  = \int_0^\infty p_0(\lp)
      \exp(-\taugeo\lp)\dd\lp
  = \mathcal L\{p_0\}(\taugeo).
  \label{eq:laplace}
\end{equation}
Thus the resolved spectrum scans the Laplace transform of a single path PDF:
wavelength changes $\taugeo$, while $p_0$ remains fixed.

\paragraph{Assumptions added by the scalar reduction.}
\begin{enumerate}
  \item Every trajectory's layer-by-layer path can be represented by one
        multiplier $\lp$.  Equivalently, scattering approximately stretches or
        shortens the direct slant path in the same proportion in all absorbing
        layers.
  \item The modeled $\taugeo(\lambda)$ contains the relevant line strength,
        pressure, temperature, and absorber-profile dependence accurately.
  \item The same $p_0(\lp)$ applies to all retained channels in the band.
        Slowly varying scattering is necessary but not sufficient: surface
        BRDF, polarization, aerosol properties, and instrumental response must
        also not introduce important channel-dependent path ensembles.
\end{enumerate}

The first assumption is the most consequential.  Molecular oxygen is well
mixed, but its line absorption is pressure and temperature dependent.  A
photon spending extra distance above a cloud and one spending it near the
surface can have the same geometrical length and different
$A_\lambda(\Gamma)$.  The general theorem, \cref{eq:general_equivalence}, still
holds; only its reduction to the scalar Laplace transform may fail.

\section{Why weak and strong channels detect different paths}

Even when a single underlying $p_0$ is exact, absorption changes the paths
represented among detected photons.  Bayes weighting of \cref{eq:laplace}
gives
\begin{equation}
  p_{\mathrm{det}}(\lp\mid\tau)
  = \frac{p_0(\lp)e^{-\tau\lp}}
         {\mathcal T(\tau)}.
  \label{eq:detected_pdf}
\end{equation}
Large $\tau$ suppresses the long-path tail.  Consequently,
\begin{align}
  -\frac{\dd\ln\mathcal T}{\dd\tau}
    &= \E_{\mathrm{det},\tau}[\lp],
    \label{eq:local_mean}\\
  \frac{\dd^2\ln\mathcal T}{\dd\tau^2}
    &= \Var_{\mathrm{det},\tau}(\lp).
    \label{eq:local_variance}
\end{align}
The local slope at a strong line is therefore the mean of an
absorption-selected, short-path population.  This does not imply an underlying
wavelength-dependent $p_0$.  A genuine violation occurs when
$P_0(\Gamma\mid\lambda)$ changes across the band, or when no common scalar
$\lp(\Gamma)$ satisfies \cref{eq:scalar_path} for all channels.

\section{Surface reflection and the measured transmittance}

For a Lambertian surface with reflectance $\gamma$ and a direct solar
normalization $F_0\mu_0$, the simplified reflected radiance is
\begin{equation}
  \pi I_\lambda
  = F_{0,\lambda}\mu_0\,\gamma\,
    \mathcal T[\taugeo(\lambda)].
  \label{eq:radiance}
\end{equation}
Define the measured normalized quantity
\begin{equation}
  T_\lambda
  \equiv \frac{\pi I_\lambda}{F_{0,\lambda}\mu_0}
  = \gamma\int_0^\infty p_0(\lp)e^{-\taugeo\lp}\dd\lp.
  \label{eq:measured_T}
\end{equation}
Then $T(0)=\gamma$.  In the project implementation,
\texttt{compute\_transmittance} evaluates the corresponding radiance/TOA-solar
normalization times $\pi$, and the fitted intercept absorbs the effective
zero-absorption reflectance.

\paragraph{Assumptions at this step.}
The Lambertian, single-reflection expression is an interpretation, not a full
description of the OCO-2 radiance.  Spectral BRDF, polarization, multiple
surface--atmosphere interactions, fluorescence, unresolved continuum
structure, or errors in the solar normalization can enter the intercept and
the fitted curvature.  Accordingly, $\exp(\text{intercept})$ should generally
be called an \emph{effective continuum reflectance}, not an exact surface
albedo.

\section{From the Laplace transform to the cumulant fit}

Let $C_n$ be the ordinary cumulants of the underlying normalized path
distribution $p_0(\lp)$.  The cumulant-generating function is
\begin{equation}
  K(t)=\ln\E[e^{t\lp}]
      =\sum_{n=1}^\infty \frac{C_n}{n!}t^n.
  \label{eq:cgf}
\end{equation}
Because the Laplace transform is the moment-generating function evaluated at
$t=-\tau$, \cref{eq:measured_T,eq:cgf} give
\begin{equation}
  \ln T(\tau)
  = \ln\gamma
    +\sum_{n=1}^\infty \frac{(-1)^n C_n}{n!}\tau^n.
  \label{eq:ordinary_cumulants}
\end{equation}

The code uses the parameterization
\begin{equation}
  \ln T(\tau)
  = b+\sum_{n=1}^{N}\frac{(-1)^n k_n}{n}\tau^n,
  \label{eq:project_fit}
\end{equation}
so the exact mapping is
\begin{equation}
  b=\ln\gamma,
  \qquad
  k_n=\frac{C_n}{(n-1)!}.
  \label{eq:scaled_cumulants}
\end{equation}
In particular,
\begin{equation}
  k_1=C_1=\E_0[\lp],
  \qquad
  k_2=C_2=\Var_0(\lp).
  \label{eq:first_two}
\end{equation}
The first two fitted coefficients are therefore the ordinary mean and variance
under the assumptions above.  For $n\geq3$, $k_n$ is a scaled cumulant rather
than the ordinary cumulant itself.  This factorial distinction matters when
comparing the project coefficients with standard probability notation.

The design matrix used by the project is consequently
\begin{equation}
  \bm X(\tau)
  = \left[1,-\tau,\frac{\tau^2}{2},-\frac{\tau^3}{3},ldots,
           \frac{(-1)^N\tau^N}{N}\right].
  \label{eq:design}
\end{equation}
The model is linear in $(b,k_1,\ldots,k_N)$ and is solved by SVD least squares,
with a bounded linear least-squares fallback if $k_1$ or $k_2$ would be
negative.  Production uses the raw, non-Savitzky--Golay coefficient set; the
parallel smoothed fit is retained as a robustness product.  The band orders are
$N=(7,3,7)$ for O$_2$ A, weak CO$_2$, and strong CO$_2$, respectively.  The
lower weak-CO$_2$ order reflects its limited optical-depth range and the poor
identifiability of high powers of $\tau$, not a different theoretical model.

\paragraph{Assumptions at the fitting step.}
\begin{enumerate}
  \item The Taylor expansion about $\tau=0$ adequately represents the retained
        finite-$\tau$ channels.  Strong line centers may lie outside the useful
        convergence region even when the equivalence theorem itself is valid.
  \item Channel errors and forward-model errors do not create structured
        residuals that the polynomial mistakes for path cumulants.
  \item $k_1$ and $k_2$ are identifiable from the available optical-depth
        range.  Positivity is physically necessary but does not guarantee that
        the complete fitted polynomial corresponds to any nonnegative PDF.
\end{enumerate}

\section{Gamma distribution as an interpretive special case}

If $p_0(\lp)$ is gamma distributed with mean $m$ and shape $\kappa$,
\begin{equation}
  p_0(\lp)
  = \frac{(\kappa/m)^\kappa}{\Gamma(\kappa)}
    \lp^{\kappa-1}\exp\!\left(-\frac{\kappa}{m}\lp\right),
  \label{eq:gamma_pdf}
\end{equation}
then
\begin{equation}
  T(\tau)
  = \gamma\left(1+\frac{m}{\kappa}\tau\right)^{-\kappa}.
  \label{eq:gamma_T}
\end{equation}
The ordinary gamma cumulants are
$C_n=(n-1)!m^n/\kappa^{n-1}$, and therefore the project coefficients obey
\begin{equation}
  k_n=\frac{m^n}{\kappa^{n-1}},
  \qquad
  k_1=m,
  \qquad
  k_2=\frac{m^2}{\kappa},
  \qquad
  \kappa=\frac{k_1^2}{k_2}.
  \label{eq:gamma_mapping}
\end{equation}
Under a strict gamma PDF, all higher coefficients are fixed by $k_1$ and $k_2$:
\begin{equation}
  k_n=k_1\left(\frac{k_2}{k_1}\right)^{n-1}.
  \label{eq:gamma_sequence}
\end{equation}

The current production fit does \emph{not} impose
\cref{eq:gamma_T,eq:gamma_sequence}.  It fits the coefficients in
\cref{eq:project_fit} independently, constraining only $k_1,k_2\geq0$.
The gamma closed form exists in \texttt{cumulant\_fit.py} but is not used by the
active estimator.  Thus $k_1^2/k_2$ is a gamma-equivalent shape diagnostic, not
evidence that the retrieved path PDF is gamma.  Flexible PDF retrievals such as
the finite-element approach of \citet{bennartz2006} illustrate why a two-
parameter gamma form can be too restrictive for multilayer or bimodal path
populations.

\section{When one wavelength-independent scalar PDF is inadequate}

The fully stratified form of the equivalence theorem remains valid if each
trajectory is described by its layer path vector
$\bm L=(L_1,\ldots,L_Z)$.  Let $p_0(\bm L)$ be its joint distribution.  Then
\begin{equation}
  T_\lambda
  = \gamma_\lambda
    \int p_0(\bm L)
    \exp\!\left[-\sum_{z=1}^{Z}\alpha_{\lambda z}L_z\right]\dd\bm L.
  \label{eq:multidimensional}
\end{equation}
This is a multidimensional Laplace transform sampled along the wavelength-
dependent vector $\bm\alpha_\lambda$.  It reduces to
\cref{eq:laplace} only when the layer path vectors are effectively collinear
with a known reference profile, $L_z\simeq\lp\,\bar L_z$.

This distinction connects directly to prior work.  \citet{scholl2006} used
pressure- and temperature-dependent O$_2$ absorption and separated path
contributions above, within, and below cloud layers.  \citet{min2004} compared
O$_2$ and water-vapor path statistics and showed that absorber vertical
inhomogeneity changes the inferred effective paths.  \citet{bennartz2006}
demonstrated the difficulty of a single classical Laplace inversion for
multilayer clouds and tested a more flexible PDF representation against Monte
Carlo path distributions.

For OCO-2, a failure of the scalar approximation could appear as systematic
differences between channels with similar total $\tau$ but different vertical
absorption kernels.  This would not disprove the equivalence theorem.  It would
show that \cref{eq:scalar_path}, rather than
\cref{eq:general_equivalence}, is too aggressive.

\section{Connection to the OCO-2 cloud-proximity project}

\begin{table}[ht]
\centering
\small
\caption{Relationship between the derivation, literature, and implementation.}
\label{tab:map}
\begin{tabular}{p{0.20\linewidth}p{0.27\linewidth}p{0.43\linewidth}}
\toprule
Concept & Literature role & Project realization and interpretation \\
\midrule
Equivalence theorem & \citet{irvine1964}, \citet{partain2000}, and
\citet{stephens2000} separate scattering-path statistics from gaseous
absorption. & The project treats channels in a narrow band as samples of one
scattering-generated path ensemble evaluated at different modeled optical
depths. \\
Narrow O$_2$ A band & \citet{heidinger2000}, \citet{funk2003}, and
\citet{min2004} exploit nearly constant continuum scattering and strongly
varying O$_2$ absorption. & O$_2$ A is fitted independently from the CO$_2$
bands.  The O$_2$ coefficients are insensitive to a real CO$_2$ enhancement
to first order and therefore provide a useful cloud/scattering fingerprint. \\
Path moments & \citet{funk2003} and \citet{scholl2006} infer PPDFs or their
first two moments from high-resolution spectra. & $k_1$ and $k_2$ are used as
mean-path and path-variance features, but should be described as
band-effective quantities under the scalar-path approximation. \\
Complex cloud paths & \citet{stephens2000} treats 3-D transport;
\citet{bennartz2006} tests flexible PDF inversion for single and multilayer
clouds. & The fitted higher-order polynomial is more flexible than a strict
gamma PDF, but it is not a full PDF inversion and cannot uniquely recover
multimodal or layer-resolved paths. \\
Operational fit & Classical inverse-Laplace retrievals are ill conditioned and
require high spectral quality \citep{min2004,bennartz2006}. & The project uses a
stable low-dimensional summary: exact linear fitting of log transmittance,
band-specific truncation, and nonnegative $k_1,k_2$. \\
Scientific role & A-band studies use path statistics to diagnose clouds and
scattering. & The spectral block is primarily the mechanistic and
plume-safety evidence.  Feature ablation shows it is approximately neutral for
TCCON predictive skill, so physical interpretation should not be overstated as
the operational source of correction skill. \\
\bottomrule
\end{tabular}
\end{table}

The project's strongest defensible statement is therefore:
\begin{quote}
Within each OCO-2 band, the fitted coefficients are effective cumulants of an
absorption-path enhancement distribution under the equivalence-theorem,
wavelength-invariant-scattering, and scalar-path approximations.  The first two
coefficients reduce to the mean and variance of a common underlying path PDF
when those approximations hold.
\end{quote}
It is stronger than calling the coefficients arbitrary polynomial features,
but more accurate than claiming an exact reconstruction of the geometrical
photon-path PDF.

\section{Manuscript-ready conservative-ensemble derivation}
\label{sec:manuscript_derivation}

The following compact formulation is recommended for the manuscript.  It
defines the fitted object unambiguously as the path distribution of the
\emph{conservative} reference problem, rather than the absorption-selected
distribution at each wavelength.

\begin{quote}
Following the equivalence theorem \citep{irvine1964,partain2000,stephens2000},
we define $P_0(\Gamma)$ as the ensemble of photon trajectories reaching the
detector in a conservative reference atmosphere with the same illumination,
viewing geometry, surface boundary, and continuum scattering properties as the
absorbing atmosphere.  Across a sufficiently narrow spectral band, these
scattering properties are assumed wavelength independent; wavelength therefore
enters only through molecular absorption.  A trajectory $\Gamma$ contributes
at wavelength $\lambda$ with Beer--Lambert weight
\begin{equation}
  w_\lambda(\Gamma)
  = \exp[-A_\lambda(\Gamma)],
  \qquad
  A_\lambda(\Gamma)
  = \int_\Gamma \alpha_\lambda(\bm r)\dd s.
  \label{eq:ms_weight}
\end{equation}
The normalized radiance is consequently
\begin{equation}
  \mathcal T_\lambda
  = \E_{P_0}[w_\lambda]
  = \int P_0(\Gamma)e^{-A_\lambda(\Gamma)}\dd\Gamma.
  \label{eq:ms_equivalence}
\end{equation}
We approximate the absorber amount along each trajectory by
$A_\lambda(\Gamma)\simeq\tau_\lambda\lp(\Gamma)$, where $\tau_\lambda$ is the
modeled direct-slant optical depth and $\lp$ is the corresponding normalized
path enhancement.  If $p_0(\lp)$ denotes the distribution of $\lp$ induced by
the conservative ensemble $P_0$, then
\begin{equation}
  \mathcal T(\tau_\lambda)
  = \int_0^\infty p_0(\lp)e^{-\tau_\lambda\lp}\dd\lp,
  \label{eq:ms_laplace}
\end{equation}
so the absorption spectrum samples the Laplace transform of one
wavelength-independent conservative path distribution.  Including an
effective continuum reflectance $\gamma$ gives
\begin{equation}
  T_\lambda
  = \frac{\pi I_\lambda}{F_{0,\lambda}\mu_0}
  = \gamma\,\mathcal T(\tau_\lambda).
  \label{eq:ms_measured}
\end{equation}
Expanding the log Laplace transform about zero absorption yields
\begin{equation}
  \ln T(\tau)
  = \ln\gamma
    +\sum_{n=1}^{\infty}\frac{(-1)^n C_n}{n!}\tau^n
  \simeq b+\sum_{n=1}^{N}\frac{(-1)^n k_n}{n}\tau^n,
  \qquad
  k_n=\frac{C_n}{(n-1)!},
  \label{eq:ms_fit}
\end{equation}
where $C_n$ is the $n$th cumulant of the conservative distribution
$p_0(\lp)$.  Hence $k_1=\E_{p_0}[\lp]$ is the mean path enhancement and
$k_2=\Var_{p_0}(\lp)$ is its variance under the stated approximations.
\end{quote}

Immediately after this derivation, the manuscript should include the following
qualification:

\begin{quote}
The conservative distribution $p_0(\lp)$ is distinct from the path
distribution among photons detected within an absorption channel.  The latter
is $p_{\mathrm{det}}(\lp\mid\tau)\propto
p_0(\lp)e^{-\tau\lp}$ and therefore varies between weak and strong channels
even when the equivalence-theorem assumptions hold.  The fitted coefficients
describe the common conservative ensemble, inferred from its wavelength-
dependent absorption weighting, rather than a separate photon-path
distribution at each wavelength.
\end{quote}

This language makes the scientific claim testable: the project assumes a
common conservative ensemble within each band, not identical detected paths at
all optical depths.  If layer-dependent absorption cannot be represented by
$A_\lambda\simeq\tau_\lambda\lp$, the coefficients remain band-effective
summaries but cease to be exact geometrical-path cumulants.

\section{Recommended validation tests}

The following tests directly target the assumptions rather than only fit
quality:
\begin{enumerate}
  \item Fit weak and intermediate optical-depth channels, then predict the
        strong channels.  Successful prediction supports one underlying
        $p_0$; structured failure identifies either truncation error or a
        scalar-path failure.
  \item Compare channels with similar total $\tau$ but different pressure or
        vertical weighting.  A persistent radiance difference after continuum
        control is evidence that scalar $\tau$ is insufficient.
  \item Use backward Monte Carlo output to compute both the true geometrical
        path histogram and layer-resolved absorber amounts.  Apply the
        production estimator to the simulated spectrum and test whether fitted
        $k_1,k_2$ close against the corresponding moments.
  \item Repeat the closure for independent-column and full 3-D transport.  This
        separates adjacency-driven horizontal paths from ordinary vertical
        cloud truncation.
  \item Examine residuals against wavelength, line identity, pressure
        broadening, surface type, and polarization after conditioning on
        $\tau$.  Dependence on wavelength beyond $\tau$ challenges the common-
        path assumption.
\end{enumerate}

\section{Assumption ledger}

\begin{table}[ht]
\centering
\small
\caption{Assumptions and consequences if they fail.}
\begin{tabular}{p{0.25\linewidth}p{0.31\linewidth}p{0.34\linewidth}}
\toprule
Assumption & Role & Consequence of failure \\
\midrule
Constant scattering properties within a band & Keeps $P_0(\Gamma)$ common to
all channels. & The spectrum is not a transform of one underlying trajectory
ensemble. \\
Absorption is passive survival weighting & Establishes the equivalence theorem,
\cref{eq:general_equivalence}. & Inelastic or coupled processes require a full
wavelength-changing transport treatment. \\
Scalar absorber path, \cref{eq:scalar_path} & Reduces the trajectory functional
to one-dimensional $p_0(\lp)$. & $k_1,k_2$ become effective absorber- and
channel-weighted coefficients rather than geometrical-path moments. \\
Accurate continuum/reference normalization & Identifies the radiance ratio with
transmittance. & Surface or instrument spectral structure leaks into the
intercept and cumulants. \\
Lambertian effective surface & Gives $b=\ln\gamma$. & The intercept is only an
effective reflectance; BRDF coupling may also affect curvature. \\
Adequate Taylor truncation & Connects finite-$\tau$ observations to derivatives
at $\tau=0$. & Strong-line information can bias or destabilize the fitted
cumulants despite a valid equivalence theorem. \\
Existence of a physical PDF & Gives cumulant meaning to all coefficients. & A
polynomial with $k_1,k_2\geq0$ may still not be the log-Laplace transform of any
nonnegative PDF. \\
Gamma PDF, only when invoked & Gives \cref{eq:gamma_mapping,eq:gamma_sequence}. &
$k_1^2/k_2$ is merely gamma-equivalent; independently fitted higher terms need
not follow the gamma sequence. \\
\bottomrule
\end{tabular}
\end{table}

\section{Conclusion}

The equivalence theorem provides the rigorous starting point: absorption
exponentially reweights trajectories generated by an otherwise unchanged
scattering problem.  The familiar one-dimensional PPDF Laplace transform
requires an additional assumption that every trajectory can be summarized by
one wavelength-independent absorber-path multiplier.  Under that assumption,
the project fit has a direct probabilistic interpretation: $k_1$ is mean path
enhancement and $k_2$ is its variance.  Weak and strong channels naturally
contain different \emph{detected} path populations, but that alone does not
violate the theorem.  The important failure mode is wavelength-dependent
scattering or layer-dependent absorption that cannot be collapsed onto one
scalar path.  The current coefficients remain scientifically useful in that
case, but should be reported as band-effective path cumulants and validated
against layer-resolved Monte Carlo path statistics.

\bibliographystyle{plainnat}
\bibliography{equivalence_theorem_fitting}

\end{document}
